% ============================================================================
%  99-ket-luan.tex — Bất đồng chưa giải quyết, khoảng trống, và việc phải làm
%  Ngày tạo: 2026-09-09 | Sửa gần nhất: 2026-09-09 | Ôn gần nhất: chưa
% ============================================================================
\documentclass[hieu-suat.tex]{subfiles}
\begin{document}
\chapter{Kết luận: cái đã chắc, cái còn cãi nhau, cái chưa ai làm}
\ptrtarget{99}

\begin{tomtat}
Chương này chốt lại ba thứ mà một khảo sát tử tế bắt buộc phải có: những gì tài liệu đã đồng thuận, những chỗ các nhóm còn mâu thuẫn thật sự, và những câu hỏi dẫn đường mà \emph{không có công trình nào} trả lời được --- tức là khoảng trống nghiên cứu, cũng là cầu nối sang thư mục \code{research/}.
\end{tomtat}

\section{Cái đã đủ chắc để làm theo}

Bốn kết luận dưới đây đến từ can thiệp có gán ngẫu nhiên, đã được lặp lại độc lập, và cỡ hiệu ứng đủ lớn để thấy được ở quy mô một người:

\begin{enumerate}
\item \textbf{Truy hồi hơn đọc lại}, và khoảng cách càng lớn khi thời gian tới lúc cần dùng càng dài \cite{roediger2006,karpicke2008,dunlosky2013}.
\item \textbf{Giãn cách và xen kẽ} tốt hơn dồn khối, kể cả khi chúng làm việc luyện tập \emph{có vẻ} kém hiệu quả ngay lúc đó \cite{cepeda2006,rohrer2015}.
\item \textbf{Sinh nội dung đối lập một cách bắt buộc} --- xét mặt đối lập, hồi tưởng từ tương lai --- làm giảm thiên lệch mạnh hơn hẳn ý chí ``cố khách quan'' \cite{lord1984,mitchell1989}.
\item \textbf{So sánh công bằng và đo nhiễu nền} là điều kiện tối thiểu để một con số trong CV/ML mang thông tin \cite{musgrave2020,melis2018,bouthillier2021}.
\end{enumerate}

\section{Ba bất đồng thật, chưa ngã ngũ}

\textbf{Một --- luyện tập có chủ đích giải thích được bao nhiêu.} Macnamara và cộng sự \cite{macnamara2014} đo được sức giải thích nhỏ và giảm dần theo lĩnh vực; Ericsson \cite{ericsson2021} đáp rằng thước đo của họ gộp cả những hoạt động không thoả định nghĩa gốc. Bằng chứng hiện có không phân giải được vì hai bên đo hai đại lượng khác nhau; muốn giải quyết cần một nghiên cứu đo riêng phần hoạt động thoả đủ năm tiêu chí, việc mà chưa ai làm ở quy mô đủ lớn \nM{}.

\textbf{Hai --- ``viết mỗi ngày''.} Boice \cite{boice1983,boice1989} có số liệu can thiệp ủng hộ mạnh; Sword \cite{sword2016} chỉ ra cỡ mẫu chín người mỗi nhóm, nhóm đối chứng nhân tạo, 15 trên 42 người bỏ cuộc, và đưa dữ liệu ngược chiều từ 100 phỏng vấn cho thấy chỉ khoảng một phần tám người thành công viết mỗi ngày. Bằng chứng hiện nghiêng về phía \emph{không có nếp viết duy nhất đúng}, nhưng chưa có thí nghiệm cỡ lớn nào thay thế \nM{}.

\textbf{Ba --- quy tắc cơ hội ngang nhau.} Mô hình của Simonton \cite{simonton1997} vẫn đang được tranh luận trong tài liệu về sáng tạo; nó dựa trên dữ liệu lịch sử và tương quan, nên hệ quả thực hành (``hoàn thành nhiều lần thử'') phải coi là giả thuyết làm việc chứ không phải kết luận.

\section{Khoảng trống: câu hỏi dẫn đường mà tài liệu không trả lời được}

Theo \code{../research-rules.md} mục 6.4, câu hỏi nào tài liệu không trả lời được thì phải nói thẳng. Có bốn câu như vậy \nM{}:

\begin{enumerate}
\item \textbf{Không có nghiên cứu can thiệp nào} đo hiệu quả của một chương trình luyện tập có cấu trúc lên năng suất hay chất lượng \emph{nghiên cứu}. Toàn bộ chương \ptr{10} vì thế là tổ hợp chưa được kiểm.
\item \textbf{Không tìm được nghiên cứu thực nghiệm} nào so sánh cách chuyên gia và người mới đọc một bài báo trong khoa học máy tính; phần đọc ở \ptr{03} phải mượn kết quả từ cờ vua và vật lý, và việc mượn đó là một giả định chưa kiểm.
\item \textbf{Không có tương đương của Dunbar cho ML/CV}: chưa ai quan sát trực tiếp một nhóm nghiên cứu học máy để đo xem suy luận phân tán hoạt động thế nào trong ngành này.
\item \textbf{Không có bằng chứng có đối chứng} cho hiệu quả của quy trình viết lại nhiều lượt; kết quả của Sommers \cite{sommers1980} là mô tả, không phải can thiệp.
\end{enumerate}

\section{Giả thuyết tôi sẽ tự kiểm}

Từ khoảng trống thứ nhất, câu hỏi hẹp nhất mà tôi trả lời được bằng chính công việc của mình \nM{}: \emph{việc bắt buộc viết dự đoán kèm xác suất trước mỗi thí nghiệm có làm giảm số thí nghiệm bị bỏ đi vì thiết kế sai không?} Giả thuyết hiện tại: có, và cơ chế là nó buộc phải nêu tiêu chí phân biệt giả thuyết trước khi tiêu tài nguyên máy --- mức tin khoảng 65\%. Điều sẽ chứng minh tôi sai: sau mười hai tuần, tỉ lệ thí nghiệm phải chạy lại vì thiết kế không đổi, hoặc tăng.

\section{Nếu chỉ làm năm việc}

\begin{enumerate}
\item Viết dự đoán kèm xác suất trước khi nhìn bất kỳ kết quả nào \cite{mellers2015}.
\item Đo nhiễu nền trước khi tuyên bố bất kỳ cải thiện nào \cite{bouthillier2021}.
\item Chạy premortem và ``xét mặt đối lập'' trước khi tin một kết quả \cite{lord1984,mitchell1989}.
\item Đưa bản thảo xấu cho một người thật, với yêu cầu bác bỏ chứ không phải góp ý \cite{mercier2011,dunbar1997}.
\item Ôn bằng truy hồi theo lịch giãn dần, không đọc lại \cite{karpicke2008,cepeda2006}.
\end{enumerate}

\end{document}
